\chapter{Why a Build System}\label{ch:mk-why}

\texttt{make} was written in 1976, is older than most of the languages it
builds, and has outlived a long procession of intended replacements. That
survival is not inertia. It solves a problem that never went away, and it
solves it with one idea small enough to fit in a paragraph.

\section{Notation for Part III}\label{sec:mk-notation}

Makefile excerpts appear in framed blocks, as shell transcripts did in
Part~\ref{part:shell}:

\begin{makecode}
program: main.o util.o        # target: prerequisites
    gcc -o program main.o util.o  # recipe
\end{makecode}

\begin{note}[Indentation in these listings]
	Recipe lines are printed indented, as above. In a real Makefile that
	indentation must be a single literal tab character --- not the spaces the page
	appears to show. \autoref{sec:mk-tab} is about nothing else, and it is the
	first thing to check when a Makefile will not run.
\end{note}

Terminal transcripts keep the convention of \autoref{sec:sh-notation}: a line
beginning \texttt{\$}~is what you type, and everything else is output.

\begin{shell}
$ make
gcc -c main.c
gcc -o program main.o util.o
\end{shell}

\begin{keys}[0.24]
	\texttt{\$(VAR)}        & A make variable being expanded. Parentheses are conventional; \texttt{\$\{VAR\}} is identical \\
	\texttt{\$\$}           & A literal dollar sign, passed through to the shell             \\
	\texttt{\{target\}}, \texttt{[prereq]} & Placeholders, as in \autoref{sec:sh-notation}  \\
\end{keys}

\begin{note}
	The \texttt{\$} character means two different things two lines apart. In
	\texttt{\$(CC)} it is \emph{make} expanding a variable; in a recipe,
	\texttt{\$\$HOME} is what you write to get the \emph{shell} a single
	\texttt{\$HOME}. Make expands first and the shell second, and forgetting the
	order is the most common Makefile bug after the tab.
\end{note}

\section{Compiling by Hand, and Where It Breaks}\label{sec:mk-byhand}

A one-file program needs no build system:

\begin{shell}
$ gcc -Wall -o hello hello.c
\end{shell}

Three files is still manageable, if tedious:

\begin{shell}
$ gcc -Wall -c main.c
$ gcc -Wall -c parser.c
$ gcc -Wall -c util.c
$ gcc -o app main.o parser.o util.o
\end{shell}

You could put those four lines in a shell script (\autoref{ch:sh-script}) and
call the problem solved. For a while it is. Then the project reaches forty
source files, a full rebuild takes ninety seconds, and you are changing one
line at a time --- at which point the script is recompiling thirty-nine files
that did not change, and you start editing the script to comment things out.
That is the moment the problem stops being ``run these commands'' and becomes
something else.

\begin{intuition}
	A build script answers \emph{what commands do I run?} A build system answers
	\emph{which of them do I still need to run?} The second question is the one
	that scales, and it cannot be answered without knowing what depends on what.
\end{intuition}

\section{Dependency Graphs and Incremental Builds}\label{sec:mk-graph}

\begin{definition}[Dependency graph]
	A directed graph whose nodes are files and whose edges say ``this file is built
	from that one''. Building a target means visiting its prerequisites first, and
	an edge is \vocab{stale} when the prerequisite is newer than the thing derived
	from it.
\end{definition}

For the three-file program above:

\begin{center}
	\begin{tikzpicture}[
			x=1mm, y=1mm, font=\footnotesize\ttfamily,
			file/.style={draw=Gray!60, rounded corners=2pt, inner sep=3pt,
					fill=Gray!5, anchor=center},
			dep/.style={-{Stealth[length=4pt]}, Gray!70!black},
			hdr/.style={dep, dashed}
		]
		\node[file] (app) at (  0,   0) {app};
		\node[file] (mo)  at (-32, -16) {main.o};
		\node[file] (po)  at (  0, -16) {parser.o};
		\node[file] (uo)  at ( 32, -16) {util.o};
		\node[file] (mc)  at (-32, -32) {main.c};
		\node[file] (pc)  at (  0, -32) {parser.c};
		\node[file] (uc)  at ( 32, -32) {util.c};
		\node[file] (h)   at (  0, -50) {parser.h};
		\draw[dep] (mo) -- (app);
		\draw[dep] (po) -- (app);
		\draw[dep] (uo) -- (app);
		\draw[dep] (mc) -- (mo);
		\draw[dep] (pc) -- (po);
		\draw[dep] (uc) -- (uo);
		% the shared header reaches two objects; route the edges through the
		% gaps between columns so neither passes behind parser.c
		\draw[hdr] (h) to[out=150, in=-35] (mo);
		\draw[hdr] (h) to[out=30,  in=-45] (po);
	\end{tikzpicture}
\end{center}

Solid edges are the obvious ones; the dashed pair is the shared header, which
both \file{main.c} and \file{parser.c} include.

Now the question ``what must I rebuild after editing \file{parser.c}?'' has a
mechanical answer: follow the arrows out of that node. \file{parser.o}, then
\file{app}. Two commands, not four. Editing \file{parser.h} instead reaches
\file{main.o} as well, because both source files include it --- and \emph{that}
edge is the one a hand-written script always forgets.

\begin{definition}[Incremental build]
	Rebuilding only the parts of the graph that are out of date. The saving is not
	a constant factor: as a project grows, the full build grows with it while a
	typical incremental build stays roughly the same size.
\end{definition}

\begin{moral}
	Everything else in Part~\ref{part:make} is syntax. The content is this graph.
	When a Makefile misbehaves --- rebuilding too much, or worse, too little ---
	the question to ask is never ``what is wrong with my rules'' but ``what does my
	graph actually say''.
\end{moral}

\section{Make's Model}\label{sec:mk-model}

\texttt{make} implements exactly the above, with one simplifying decision that
explains most of its behaviour and all of its failure modes.

\begin{definition}[Make's rule]
	Given a \vocab{target}, a list of \vocab{prerequisites}, and a \vocab{recipe}
	of shell commands: build each prerequisite first, then run the recipe if and
	only if the target does not exist, or any prerequisite has a
	\emph{modification time} newer than the target's.
\end{definition}

\begin{intuition}
	Make does not know what your compiler does, what a \texttt{.c} file is, or what
	changed inside it. It compares timestamps on files. That is the whole of its
	understanding, and it is why make works equally well for C, \LaTeX{}, image
	conversion and data pipelines --- and why it is fooled by anything that does
	not show up as a timestamp.
\end{intuition}

\begin{note}[What timestamps cannot see]
	Three consequences follow immediately, and all three will bite you eventually:
	\begin{enumerate}[(i)]
		\item Touching a file without changing it (\cmd{touch main.c}) forces a rebuild
		      that achieves nothing.
		\item Changing a \emph{compiler flag} rebuilds nothing at all, because no file
		      changed --- your object files are now inconsistent and make is content.
		\item A clock skew, or a file copied with its old timestamp, can make a
		      prerequisite look older than its target, and the rebuild silently does not
		      happen.
	\end{enumerate}
	\autoref{sec:mk-depgen} fixes the first of these properly. For the second the
	answer is usually \cmd{make clean}.
\end{note}

\section{Where Make Sits}\label{sec:mk-alt}

\begin{keys}[0.24]
	\texttt{make}        & Rules and timestamps, by hand. Universal, tiny, and unopinionated \\
	\texttt{CMake}       & Generates Makefiles (or Ninja files) from a higher-level description; handles cross-platform and dependency discovery \\
	\texttt{Ninja}       & A make-like executor designed to be generated rather than written; much faster on large graphs \\
	\texttt{latexmk}     & A purpose-built builder for \LaTeX{} that knows about reruns, \texttt{biber} and \texttt{makeindex} (\autoref{sec:setup-vimtex}) \\
	\texttt{cargo}, \texttt{go build} & Language-specific; the dependency graph is derived from the source, not declared \\
\end{keys}

\begin{remark}
	CMake and Ninja have not replaced make so much as moved it. CMake's output is
	usually a Makefile, and reading it is the fastest way to appreciate why nobody
	writes those by hand. The reason to learn make anyway is that it is present on
	every Unix machine, needs no project, and is the right size for the jobs that
	are too small to justify a real build system --- which is most of them.
\end{remark}

\begin{moral}
	Learn \texttt{make} for your own projects and for reading other people's. Reach
	for CMake when you need to build on Windows too, or when finding libraries
	becomes a project of its own.
\end{moral}
